\chapter*{Introduction générale}
\markboth{Introduction générale}{}
\addcontentsline{toc}{chapter}{Introduction générale}
L'accès à un stage ou à une première expérience professionnelle met en relation deux besoins complémentaires. Les étudiants cherchent des missions compatibles avec leur formation, leurs projets et les compétences qu'ils souhaitent approfondir. Les entreprises souhaitent identifier des profils capables de contribuer à leurs activités, tout en tenant compte du caractère formateur du stage. La rencontre entre ces besoins dépend de la qualité des informations disponibles et de la manière dont elles sont comparées.

Les curriculum vitæ et les descriptions d'offres constituent des documents riches, mais peu homogènes. Une technologie peut être citée par son nom, par un framework associé ou par une expression propre à une langue. Les expériences sont décrites selon des formats différents et les compétences déclarées n'ont pas toutes le même niveau de preuve. Une plateforme qui ne fait que stocker ces documents laisse à l'utilisateur l'intégralité du travail de rapprochement. À l'inverse, une plateforme qui réduit ce travail à une note opaque peut donner une impression de précision sans fournir les éléments permettant d'apprécier le résultat.

Ce mémoire étudie la problématique suivante : \emph{comment concevoir une plateforme de mise en relation étudiants-entreprises qui automatise une partie de la comparaison entre profils et offres, tout en conservant des résultats explicables, une architecture évolutive et des limites identifiables ?} La réponse proposée prend la forme de SmartRecruit, développé dans le cadre du stage MDL/26/12 réalisé chez Best Solutions. Le projet associe la gestion des parcours de recrutement à un moteur de recommandation fondé sur des représentations textuelles et des connaissances explicites sur les compétences.

Le premier objectif est fonctionnel : regrouper les profils étudiants, les offres publiées, les candidatures et les décisions dans une application cohérente. Le deuxième objectif est technique : séparer les responsabilités de l'application web, de la persistance et du traitement de texte, de manière à pouvoir faire évoluer le moteur sans réécrire tous les parcours. Le troisième objectif est méthodologique : organiser le travail en flux avec Kanban, relier les besoins aux modèles UML et définir des critères vérifiables pour chaque élément réalisé.

Kanban fournit le cadre de conduite retenu pour cette présentation du projet. Les demandes sont décrites par des cartes de travail, priorisées dans une réserve puis tirées vers la réalisation selon la capacité disponible. La limitation du travail en cours évite de multiplier les fonctionnalités partiellement terminées. La vérification des droits, des données et des interfaces intervient avant le passage d'une carte à l'état terminé. L'organisation présentée est reconstituée à partir du périmètre logiciel disponible ; elle ne suppose pas l'existence d'un historique de tableau ou de mesures de flux établis pendant le stage.

La conception UML rend explicites les responsabilités et les interactions : cas d'utilisation par rôle, structure des entités métier, échanges lors d'une candidature ou d'une recommandation, activité de traitement du CV et états de suivi. Les modèles servent de lien entre les exigences et les opérations effectivement codées. La partie réalisation illustre ce lien par des captures du rendu des interfaces Blade actuelles, accompagnées de descriptions centrées sur les actions de l'utilisateur.

Le choix initial s'oriente vers une méthode légère combinant TF-IDF, similarité cosinus et couverture de compétences. Cette méthode constitue une base de comparaison contrôlable : elle permet d'inspecter les termes, les synonymes et les pondérations utilisés. Elle ne résout cependant pas toutes les difficultés linguistiques. La négation, le niveau de maîtrise, les compétences seulement évoquées et les documents numérisés comme images restent des sources d'erreur. L'architecture prévoit donc une possibilité d'évolution, sans confondre cette possibilité avec une intégration déjà réalisée.

L'apport du travail ne se limite pas au calcul d'une note. Il concerne aussi l'articulation entre cette note et les processus applicatifs : à quel moment la calculer, quelles données lui transmettre, comment la conserver, et comment réagir lorsque le service chargé du calcul devient indisponible. Le repli de calcul PHP doit être distingué de la persistance : MySQL est le magasin métier utilisé par les routes actuelles et son indisponibilité entraîne une erreur explicite. L'ancien magasin de démonstration JSON n'est plus sélectionné comme repli par ces routes.

Le mémoire est organisé en six chapitres. Le chapitre~\ref{chap:contexte} présente le contexte, la problématique, l'état de l'art et la démarche Kanban. Le chapitre~\ref{chap:besoins} formalise les exigences et leur traduction en cartes de travail assorties de critères d'acceptation. Le chapitre~\ref{chap:conception} expose la conception UML et les choix d'architecture. Le chapitre~\ref{chap:realisation} présente l'implémentation et les interfaces web illustrées. Le chapitre~\ref{chap:moteur} détaille les traitements et les formules de recommandation. Le chapitre~\ref{chap:validation} distingue les vérifications du rendu actuel, les sondes historiques et les essais restant à conduire. Les annexes rassemblent le dictionnaire de données, les contrats d'échange et les éléments de reproductibilité.
